\section{The complete transfer from disjoint coset systems}
\label{l:core-section}

The coset bounds are now available on one common measure. We first turn them into an operator estimate by following the orbit of one vector. For eight complementary coset systems, this gives convergence after four convolutions. Central projections then retain the whole Fourier matrix and reduce that count to two. Both arguments keep every restriction multiplicity; they differ in what they average.

\begin{lemma}[The orbit of a single vector]\label{l:single-orbit}
Let $H$ act unitarily on $V$, and suppose a vector $v$ has components only in $H$-types belonging to a set $\Lambda\subseteq\Irr(H)$. Then
\[
 \dim\Span\{\rho(h)v:h\in H\}\le\sum_{\tau\in\Lambda}(\dim\tau)^2.
\]
All copies of a type are included in this estimate.
\end{lemma}
\begin{proof}
Regroup the entire $\tau$-isotypic space as $V_\tau\otimes\mathcal A_\tau$, where $H$ acts as $\tau\otimes I$. If $v_\tau$ is the component of $v$ there, its orbit lies in the image of
\[
 \operatorname{End}(V_\tau)\longrightarrow V_\tau\otimes\mathcal A_\tau,
 \qquad Q\longmapsto(Q\otimes I)v_\tau.
\]
The image has dimension at most $(\dim\tau)^2$, independently of $\dim\mathcal A_\tau$. Equivalently, writing $v_\tau=\sum_{j=1}^{\dim\tau}e_j\otimes a_j$, the orbit lies in $V_\tau\otimes\Span(a_j)$. Sum over inequivalent types.
\end{proof}

Write $a=\dim\tau$, $W_\tau=\Span\{\rho(h)v_\tau:h\in H\}$, and $T(Q)=(Q\otimes I)v_\tau$. The key point is the map in Figure~\ref{l:orbit-figure}: multiplicity enlarges the ambient space, but not the matrix space that moves a fixed vector.

\begin{figure}[ht]
\centering
\begin{tikzpicture}[node distance=15mm,>=Stealth,
 box/.style={draw,rounded corners,align=center,inner sep=6pt}]
 \node[box] (mat) {$\operatorname{End}(V_\tau)$\\dimension $a^2$};
 \node[box,right=22mm of mat] (orbit) {$\operatorname{im}T\supseteq W_\tau$\\dimension at most $a^2$};
 \node[box,below=10mm of orbit] (ambient) {$V_\tau\otimes\mathcal A_\tau$\\arbitrary multiplicity};
 \draw[->] (mat)--node[above,font=\small] {$T$} (orbit);
 \draw[->] (orbit)--node[right,font=\small] {inclusion} (ambient);
\end{tikzpicture}
\caption{The orbit of one vector in all copies of a subgroup type of degree $a$. The restriction of the group action factors through a matrix space of dimension $a^2$.}
\label{l:orbit-figure}
\end{figure}

\begin{proposition}[Orbit-span transfer]\label{l:orbit-span}
Under the assumptions of Theorem~\ref{l:main}, for every $u>0$,
\[
 \|\widehat f(\rho)\|_{\op}\le\sqrt{bBC_u}\,D^{-1/2+(u+2)/(2b)}.
\]
\end{proposition}
\begin{proof}
The disjoint subgroups form a direct product $H=H_1\times\cdots\times H_b$. Its isotypic decomposition in $V_\rho$ is an orthogonal sum of spaces
\[
 (V_{\tau_1}\otimes\cdots\otimes V_{\tau_b})\otimes\mathcal A_{\boldsymbol\tau}.
\]
Every occurring tensor product has dimension $\prod_i\dim\tau_i\le D$. Assign its whole isotypic space to the least index $i$ for which $\dim\tau_i\le D^{1/b}$. Group the assigned spaces to obtain $V_\rho=E_1\oplus\cdots\oplus E_b$. Equivalently, let $P_i$ select the small $H_i$-types. These projections commute and $\prod_i(I-P_i)=0$. The projections $Q_i=P_i\prod_{j<i}(I-P_j)$ are orthogonal and sum to identity: on each product type $Q_i$ is identity exactly when $i$ is its first small factor. Thus the explicit splitting $v_i=Q_iv$ is precisely the least-index assignment above.

Fix $v_i\in E_i$ and put $W_i=\Span\rho(H_i)v_i$. Lemma~\ref{l:single-orbit} and the definition of $C_u$ give
\begin{equation}\label{l:orbit-dimension}
 \dim W_i\le\sum_{\tau\in\Irr(H_i):\,\dim\tau\le D^{1/b}}(\dim\tau)^2
 \le C_uD^{(u+2)/b}.
\end{equation}
The second inequality follows term by term from $(\dim\tau)^2\le D^{(u+2)/b}(\dim\tau)^{-u}$. It is uniform in the block size.

The remaining task is to use this dimension bound on a coefficient. Let $\Pi_i$ be the orthogonal projection onto $W_i$. Since $W_i$ is $H_i$-invariant, the projection onto $\rho(g)W_i$ depends only on $gH_i$. The coset cap therefore gives, for every $w$,
\begin{align*}
 \sum_g f(g)\|\Pi_{\rho(g)W_i}w\|^2
 &\le B\E_{g\sim U_G}\|\Pi_{\rho(g)W_i}w\|^2\\
 &=B\frac{\dim W_i}{D}\|w\|^2.
\end{align*}
Indeed $\E_U\rho(g)\Pi_i\rho(g)^*$ commutes with $\rho(G)$ and has trace $\dim W_i$, so Schur's lemma identifies it as $(\dim W_i/D)I$.

Since $\rho(g)v_i\in\rho(g)W_i$ and $f$ has mass at most one, Cauchy--Schwarz yields
\[
 |\langle w,\widehat f(\rho)v_i\rangle|
 \le\|v_i\|\left(\sum_gf(g)\|\Pi_{\rho(g)W_i}w\|^2\right)^{1/2}
 \le\sqrt{BC_u}D^{-1/2+(u+2)/(2b)}\|w\|\|v_i\|.
\]
For $v=\sum_i v_i$, orthogonality gives $\sum_i\|v_i\|\le\sqrt b\|v\|$. Summing proves the claim. In particular no dimension of a multiplicity space has been discarded.
\end{proof}

\begin{corollary}[Four convolutions by the orbit argument]\label{l:orbit-completion}\label{l:eight-convolutions}
Let $n\to\infty$ through multiples of eight. Suppose probabilities $\mu_n$ on $S_n$ have sign mean $o(1)$ and
\[
 \E_{|M|=n/8}\|\mu_{n,H_M}-U_{G/H_M}\|_{\TV}=o(1).
\]
Then $\|\mu_n^{*4}-U_G\|_{\TV}\to0$.
\end{corollary}
\begin{proof}
Lemma~\ref{l:partition-selection} selects an eight-block partition with summed marginal error $o(1)$. Lemma~\ref{l:trim} supplies a common subprobability $f_n\le\mu_n$ of mass $1-o(1)$ and coset cap one. Proposition~\ref{l:orbit-span} at $u=1$ gives $\|\widehat f_n(\rho)\|_{\op}\le\sqrt{8C_1}D^{-5/16}$. The nonsign Plancherel contribution after four factors is at most $(8C_1)^4 Z_{1/2}(n)$, because $2-8(5/16)=-1/2$. Sign and missing mass tend to zero by Proposition~\ref{l:amplification}. This completes a route from averaged observations to the full law.
\end{proof}

\subsection{Central projections retain the whole Fourier matrix}

The orbit argument controls one vector at a time. To estimate the
Hilbert--Schmidt norm without paying an extra ambient dimension, we
instead convolve with a central idempotent of the subgroup. Its
cosetwise squared norm controls all Fourier matrix entries at once.

For an irreducible $H$-type $\tau$ of degree $a$, define the mass function
\[
 e_{H,\tau}(h)=\frac{a}{|H|}\overline{\chi_\tau(h)}\quad(h\in H),\qquad e_{H,\tau}=0\text{ off }H.
\]
Character orthogonality shows that $\widehat e_{H,\tau}(\rho)$ is the orthogonal projection $P_{\rho,H,\tau}$ onto the full $\tau$-isotypic subspace in $\rho|_H$. Thus every copy is included.

\begin{lemma}[Cosetwise isotypic estimate]\label{l:central-isotypic}
Let $H\le G$ and let $f\ge0$ have mass $z\le1$ and satisfy $f(gH)\le B/[G:H]$. For an irreducible $H$-type $\tau$ of degree $a$,
\begin{equation}\label{l:central-global}
 \sum_{\rho\in\Irr(G)}D_\rho\|\widehat f(\rho)P_{\rho,H,\tau}\|_{\HS}^2\le Bza^2.
\end{equation}
\end{lemma}
\begin{proof}
Convolution with $e=e_{H,\tau}$ acts separately in each left coset $Q=xH$. Since $\sum_{h\in H}|e(h)|^2=a^2/|H|$, the triangle inequality for a positive weighted sum of translates gives
\[
 \sum_{h\in H}|(f*e)(xh)|^2\le f(Q)^2\frac{a^2}{|H|}.
\]
One may equivalently normalize the weights by $f(Q)$ and use Jensen; the zero-mass coset contributes zero. Also $\sum_Q f(Q)^2\le Bz/[G:H]$. Plancherel on $G$ proves \eqref{l:central-global}.

There is a second exact expression for the same calculation. Let $f_Q(h)=f(xh)$ and $\widehat f_Q(\tau)=\sum_hf_Q(h)\tau(h)$. Subgroup Plancherel yields
\[
 \sum_\rho D_\rho\|\widehat f(\rho)P_{\rho,H,\tau}\|_{\HS}^2
 =[G:H]\sum_Q a\|\widehat f_Q(\tau)\|_{\HS}^2.
\]
Unitarity gives $\|\widehat f_Q(\tau)\|_{\HS}\le\sqrt a\,f(Q)$ and recovers the same bound. This identity makes the normalization of the Hilbert--Schmidt norm explicit.

There is also a positive-projection proof of the same bound.
Right convolution by $e$ is an orthogonal projection, since $e*e=e$ and $e^*(h)=\overline{e(h^{-1})}=e(h)$. With the counting inner product,
\begin{align}
 |G|\|f*e\|_2^2
 &=|G|\langle f,f*e\rangle\notag\\
 &\le a^2|G|\langle f,f*U_H\rangle
 \le Bza^2.\label{l:isotypic-positive-projection}
\end{align}
Indeed $|e(h)|\le a^2U_H(h)$, so positivity of $f$ bounds the absolute convolution pointwise; and $(f*U_H)(g)=f(gH)/|H|\le B/|G|$. This proves the estimate from idempotence and positivity, independently of the cosetwise norm calculation.
\end{proof}

\begin{proposition}[Isotypic transfer]\label{l:isotypic}
Under the assumptions of Theorem~\ref{l:main},
\[
 \|\widehat f(\rho)\|_{\HS}^2\le bBC_uD^{-1+(u+2)/b}.
\]
In fact the right side can be multiplied by $f(G)$.
\end{proposition}
\begin{proof}
For each $i$, take all $H_i$-isotypic projections with degree at most $D^{1/b}$. On every product type in the restriction to $\prod_iH_i$, at least one such projection is identity. The projections commute and include all multiplicities, so
\begin{equation}\label{l:projection-cover}
 I\preceq\sum_{i=1}^b\sum_{\tau:\dim\tau\le D^{1/b}}P_{\rho,H_i,\tau}.
\end{equation}
Trace this inequality against the positive operator $\widehat f(\rho)^*\widehat f(\rho)$, apply Lemma~\ref{l:central-isotypic} after discarding other ambient irreducibles, and use \eqref{l:orbit-dimension}. The resulting inequality is $D\|\widehat f(\rho)\|_{\HS}^2\le Bf(G)bC_uD^{(u+2)/b}$.
\end{proof}

This proves Theorem~\ref{l:main}. The improvement from operator to Hilbert--Schmidt norm changes the convolution count because the regular-representation weight is then $D$, rather than the $D^2$ obtained by bounding $\|A\|_{\HS}^2$ by $D\|A\|_{\op}^2$.

\begin{corollary}[Two convolutions from eight complements]\label{l:isotypic-completion}\label{l:four-convolutions}
Under the hypotheses of Corollary~\ref{l:orbit-completion}, $\|\mu_n^{*2}-U_G\|_{\TV}\to0$.
\end{corollary}
\begin{proof}
Use the same partition and subprobability as in that corollary. Proposition~\ref{l:isotypic} with $u=1$ gives HS squared at most $8C_1D^{-5/8}$. Two powers have nonsign contribution at most $(8C_1)^2 Z_{1/4}(n)$, since $1-2(5/8)=-1/4$. Proposition~\ref{l:amplification} handles sign and mass.
\end{proof}

The last corollary is an abstract consequence of the stated marginal hypotheses. A shuffle application must supply those hypotheses at its own time and preserve its own time convention; a numerical mixing statement is not part of this abstract deduction.
